% ============================================================
%  chapters/00_abstract.tex  —  中英文摘要
% ============================================================

% ===== 中文摘要 =====
\begin{center}
{\zihao{3}\heiti \husttitle}
\end{center}

\vskip 14pt

\begin{center}
{\zihao{3}\heiti 摘\ \ 要}
\end{center}

\vskip 14pt

{\zihao{-4}\songti
脑卒中是造成长期残疾的第一大原因，在全世界范围内它也是主要的致残因素。大约有60\%-73\%的人群在经历脑梗塞之后会伴随有不同程度的视力障碍，即同向性偏盲、眼球转动受限以及单侧空间忽略等状况出现，并且这种视觉上的问题会对患者的生活自理能力和康复效果产生较大的影响[1-4]。但是现有的康复手段存在着“两端不足”的情况，传统的人工训练没有统一的标准和量化的过程记录，高端的VR、AR等设备价格昂贵、使用不便，不能适应家庭长时间康复的需求。为此，本文设计并实现了基于STM32H723VGT6微控制器的一种低成本头戴式眼部康复训练装置。

本装置使用DM-MC02核心板作为硬件平台，包括BMI088六轴惯性传感器、SG90舵机、650nm红色激光模块、CI1302语音交互模块和WS2812状态指示灯等元件。软件上采用四元数融合的方法来实现各种姿态估计算法，在没有使用任何预设的姿态参数的情况下，可以得到佩戴者的头部位置变化情况。开发出一个包含九种状态的状态机，训练模式有完整的训练模式A-E自成一体以及自选训练模式语音命令一个个的选择。五个临床训练模式，即注视稳定训练（15秒），扫视训练（8次随机跳变），平稳追踪训练（8点平滑移动），视觉聚焦训练（5循环远近交替），空间忽略训练（6次左右交替）都是从被动注视到主动视觉搜寻的一个逐渐提高的过程。语音交互系统依靠离线识别引擎，具备10条能够被识别的命令词和46条TTS播报词，并且配有三级安全防护措施，即当出现姿态超出界限时会自动停止，语音紧急中断并且能够跨过不同状态之间的转换。

测试结果表明，系统各模块通信正常，姿态估计精度满足训练需求，五种训练模式均能按照设计逻辑正确执行。本装置为脑卒中后视觉障碍的居家康复训练提供了一套功能完整、操作便捷的工程技术方案。
}

\vskip 14pt

\noindent
{\zihao{-4}\heiti 关键词：}{\zihao{-4}\songti 脑卒中；眼部康复；STM32；IMU姿态解算；语音交互}

\clearpage

% ===== 英文摘要 =====
\begin{center}
{\zihao{3}\textbf{Design and Implementation of a Head-Mounted \\[2pt]
Ocular Rehabilitation Device for Stroke Patients}}
\end{center}

\vskip 14pt

\begin{center}
{\zihao{3}\textbf{ABSTRACT}}
\end{center}

\vskip 14pt

{\zihao{-4}
Stroke is one of the leading causes of long-term disability worldwide. Between 60\% and 73\% of stroke survivors exhibit varying degrees of visual dysfunction, including homonymous hemianopsia, oculomotor disorders, and unilateral spatial neglect, all of which profoundly impair daily living activities and rehabilitation outcomes~\cite{rowe2021ivis,liu2023vision,zheng2025stroke,sun2024pursuit}. Yet current clinical interventions face a critical dichotomy: conventional manual training methods lack standardization and quantitative documentation, while high-end VR/AR-based devices remain prohibitively expensive and operationally complex for home-based long-term rehabilitation. Addressing this gap, this thesis presents the design and implementation of a low-cost, head-mounted ocular rehabilitation training device centered on the STM32H723VGT6 microcontroller.

The device employs a DM-MC02 core board as its hardware platform, integrating a BMI088 six-axis inertial measurement unit, an SG90 dual-axis servo gimbal, a 650\,nm laser module, a CI1302 AI voice interaction module, and a WS2812 programmable RGB status indicator. On the software side, a quaternion-based attitude estimation algorithm is implemented, alongside a direct mapping from the sensor coordinate frame to the head coordinate frame and a 100-frame sliding-window variance-based head stability index, achieving real-time head posture analysis at a 10\,ms cycle rate. A nine-state system state machine has been developed, supporting two operating modes: a Full Mode that automatically chains all five training patterns in sequence, and a Custom Mode in which the patient selects individual patterns via voice commands. Five clinical training protocols — fixation stability training, saccade training , smooth pursuit training , visual accommodation training , and spatial neglect training  — form a progressive rehabilitation framework spanning from passive gaze maintenance to active visual exploration. The voice interaction subsystem, built on the CI1302 offline recognition engine, supports ten recognizable command words and 46 TTS broadcast phrases, and together with a three-tier safety protection mechanism ensures training continuity and operational safety.

Test results demonstrate normal communication across all hardware modules, attitude estimation accuracy sufficient for training requirements, and correct execution logic of all five training protocols. The device provides a functionally complete and user-friendly home-based visual rehabilitation solution, offering a practical engineering pathway for long-term recovery training for post-stroke visual impairment.

}

\vskip 14pt

\noindent
{\zihao{-4}\textbf{KEY WORDS：}}{\zihao{-4} Stroke; Ocular Rehabilitation; STM32; IMU Attitude Estimation;Voice Interaction}
\clearpage
